% Pursuing Stacks (cotes 134-2 to 134-8): archive cote 134-7, site page 19 (PDF page 20), typescript page 461.
% typescript. Draft compared: ps6.tex:884-942.
% AI-transcribed by Claude Opus 5.5 (claude-opus-5-5, adaptive thinking, effort high), Message Batches API, batch ps3 (msgbatch_01XdFmuVVBAcgc7EFxdtCqsj), 2026-11-12; not yet checked by eye.
\scanpage{134-7/19}
\mspage{461}
dévissage of a given homotopy type,  and replacing every homotopy group \hand{$\pi$}$_i$
by a shifted projective resolution (a two-step resolution will do here)
M(i)$_{\hand{$\bullet$}}$\hand{,}  as indicated on page 456,  and using the corresponding semisimpli-%
cial $\underline{Z}$-modules  M(i)$_*$  as fibers,  in the successive fiberings.  To see
that what is being done on the ``discrete'' level,  working with semisimpli-%
cial \uline{sets},  can be``followed'' in an essentially unique way on the ``schema-%
tic'' level,  we hit now of course upon the key-difficulty,  pointed out
on page 444,  about the Postnikov cohomology group $H^{n+2}(X_{\deleted{n+1}\hand{(n+1)}},\hand{$\pi$}_{n+1})$
being isomorphic to the corresponding ``schematic'' group.  Now that this is
so indeed should follow from the homology isomorphisms (w) I hope,
\xout{byxpassingxoverxtoxthexdual} just dualizing the result\inserted{to cohomology}.  All this seems to
sound kind of reasonable,  it seems even that for a given homotopy type
in Hot\hand{$_0$}\deleted{\unsure{(W)}},  we should be able to squeeze out this way a unique \uline{isomorphism
type},  at any rate,  of semisimplicial unipotent bundles  -  but to see
whether it does work,  or if there is some ma\hand{j}or blunder which turns
the whole into nonsense,  will come out only from careful, down-to-earth
work,  which I am not prepared to dive into.

It occured to me that the ``Kan-Postnikov'' complexes in U(k) have
some special features among all possible complexes with $X_0 = e$,  and also
that some extra features \uline{are}  needed,  if we want the maps in (w) to be
quasi-isomorphisms.  I want to dwell upon this a little.  First of all,
the condition  $X_0$=e is indeed essential,  as we see by taking a constant
complex with value $X_0$ ,  then the homology of the three\inserted{chain} complexes
(w)  reduces to degree zero,  and the $H_0$'s are respectively
\[ \underline{Z}^{\hand{(}X_0(\underline{Z})\hand{)}}\inserted{$=\underline{Z}^{(M_0)}$} ,\quad L(X_0) = \Gamma_{\underline{Z}}{}^{\wedge}(M_0) ,\quad L_{pt}(X_0) = \Gamma_{\underline{Z}}(M_{\unsure{0}}) , \]
where $M_0$ is the $\underline{Z}$-module giving rise to $X_0$  -  and \deleted{\unsure{all}}\inserted{none of} the two maps \deleted{are}\hand{is}
\deleted{non}\inserted{an} isomorphis\over{ms}{m}\hand{,} unless $M_0$=0 .

Take now the next simple\inserted{st} case,  when $X_*$  comes from a monoid object\inserted{G}
in U(k) in the usual way;  then what we are after,  in \inserted{dual} terms of cohomology
rather than homology (taking the dual complexes of those in (w)), amounts
essentially to asking whether the usual  discrete cohomology of the
discrete  monoid  G($\underline{Z}$)  can be computed,  using \uline{polynomial} \uline{cochains}
rather than arbitrary ones.  Now this we did admit as ``well known'' in
the most evident case of all,  when G being represented \deleted{by} as an object
of U($\underline{Z}$)  by a \inserted{projective (hence free)} $\underline{Z}$-module M ,  the multiplication law is just \deleted{the} usual addi-%
tion.  It still looks reasonable enough when the monoid \deleted{law}\inserted{G} is a \uline{group} ,
with M of finite type say.  In this case,  the \inserted{Borel} theory of algebraic
\margin{19}
